\chapter{Pregnancy Warning Signs}
\index{Pregnancy Warning Signs}

\medskip
\noindent\textit{\textbf{Under review.} This chapter is an AI-assisted educational draft; it is not a substitute for professional medical advice.}

\section*{Definition and Overview}
Pregnancy warning signs are symptoms during pregnancy that may indicate a complication and need prompt assessment. The correct response depends on gestational age, pregnancy history, medical conditions, and local access to maternity care. When unsure, contact the maternity unit, midwife, obstetric clinician, or other qualified healthcare professional caring for the pregnancy.

\section*{Clinical Features}
Urgent warning signs include vaginal bleeding, severe or persistent abdominal pain, leaking fluid, regular painful contractions before term, severe headache, visual changes, sudden swelling of the face or hands, fever, painful urination with fever, persistent vomiting with inability to keep fluids down, reduced or changed fetal movement after movements have become established, chest pain, breathlessness, fainting, or one-sided leg swelling and pain.

\section*{When to Seek Medical Care}
Contact the maternity service or clinician promptly for any new, persistent, or worrying symptom. Contact your local emergency services immediately for heavy bleeding, severe abdominal pain, collapse, seizures, severe breathlessness, chest pain, confusion, or a possible stroke. Do not wait for a routine appointment when symptoms are severe or rapidly worsening.

\section*{Diagnosis and Investigations}
Assessment may include vital signs, abdominal examination, urine testing, blood tests, ultrasound, fetal heart-rate assessment, and other investigations selected for the gestational age and symptoms. Some tests and treatments are time-sensitive, so assessment should not be delayed while trying to decide which diagnosis is most likely.

\section*{Management}
Management depends on the cause and may include observation, fluids, medicines that are appropriate in pregnancy, treatment of infection, monitoring of the pregnancy, or urgent obstetric care. Do not start, stop, or change medicines, supplements, or recreational substances without checking with a qualified clinician or pharmacist.

\section*{Complications}
Possible complications include miscarriage, ectopic pregnancy, pre-eclampsia, premature labour, placental problems, infection, blood clots, and maternal or fetal compromise. Early assessment can reduce harm, but symptoms cannot be interpreted reliably without considering the stage of pregnancy and the person's clinical context.

\section*{Prognosis and Self-Care}
Many pregnancy symptoms have benign causes, but the outlook depends on the diagnosis and how quickly complications are treated. Keep scheduled antenatal visits, follow the agreed plan for monitoring fetal movement and blood pressure, stay hydrated when able, and seek help again if symptoms persist, recur, or worsen.

\section*{Sources and Further Reading}
\begin{itemize}
    \item World Health Organization. \textit{Maternal health.} https://www.who.int/health-topics/maternal-health
    \item World Health Organization. \textit{Antenatal care.} https://www.who.int/health-topics/antenatal-care
    \item World Health Organization. \textit{Pregnancy, childbirth, postpartum and newborn care: a guide for essential practice.} https://www.who.int/publications/i/item/9789241547260
\end{itemize}
